\chapter{Binary Outcomes and Maximum Likelihood}\label{ch:b3:binary-outcomes}

Many of the outcomes economists care about are yes or no. A person works or
does not; a household buys a car or does not; a commuter takes the bus or
drives. The regression methods of the earlier chapters can be applied to a
variable that takes the values 0 and 1, and often that is the right thing to
do. But a probability cannot exceed one, the effect of a variable on a
probability cannot be the same at 2~\% as at 98~\%, and a choice among several
options has a structure that a regression ignores. This chapter builds the
models that respect those facts, the method that estimates them --- maximum
likelihood --- and the model of choice among options that economists use to
forecast demand for products that do not yet exist.

\begin{openingchart}{Share of men with any earnings in 1978 against their
earnings in 1975, in thousands of dollars (15{,}992 men of a 1970s United
States population survey, the comparison sample of the job-training study of
\cref{ch:b3:difference-in-differences}; dots are shares in bins, the first one
the 1{,}748 men with no earnings in 1975). Lines: three models of the probability
on 1975 earnings alone. Earnings above about \usd{25000} are
top-coded.}{Dehejia and Wahba's NSW files (survey comparison sample); our
computation}
\begin{tikzpicture}
\begin{axis}[omscatter, width=0.78\linewidth, height=0.45\linewidth, xmin=-0.5, xmax=26.5,
  ymin=0.4, ymax=1.05, xlabel={Earnings in 1975 (thousands of dollars)},
  ylabel={Share with earnings in 1978}, legend pos=south east,
  y tick label style={/pgf/number format/.cd, fixed, precision=1}]
  \addplot+[mark size=2pt] table[x=earn75, y=share] {figdata/bachelor-3/05-binary-outcomes/out/work_bins.csv};
  \addplot[omDef, thick, no marks, sharp plot] table[x=earn75, y=lpm] {figdata/bachelor-3/05-binary-outcomes/out/work_fits.csv};
  \addplot[omThm, thick, no marks, sharp plot] table[x=earn75, y=logit] {figdata/bachelor-3/05-binary-outcomes/out/work_fits.csv};
  \addplot[omMeth, thick, dashed, no marks, sharp plot] table[x=earn75, y=probit] {figdata/bachelor-3/05-binary-outcomes/out/work_fits.csv};
  \draw[guide] (axis cs:-0.5,1) -- (axis cs:26.5,1);
  \legend{share in the bin, linear, logit, probit}
\end{axis}
\end{tikzpicture}
\end{openingchart}

Read the chart from the left. Of the men with no earnings in 1975, 44~\% had some
in 1978; of those who earned a few thousand dollars, 85--87~\%; of those near
the top, about 95~\%. % figdata: work_bins
The straight line rises at the same rate everywhere and passes 1 at the top
of the range. The logit and probit curves bend toward 1 and are hard to tell
apart. And all three miss the first dot badly: they predict 66--70~\% for men
with no earnings, % figdata: work_fits
against an observed 44~\%. Men with no earnings at all are not simply men with
low earnings, and no choice of curve fixes a model that treats them as if they
were (\cref{exo:b3:binary-outcomes:5}). The chapter's tools are worth
learning, but the chart's first lesson is that the specification matters more
than the curve.

\begin{recall}
The Year 2 volume estimated regressions by least squares, with
heteroskedasticity-robust standard errors. Logistic regression as a tool for
classifying is treated in the artificial-intelligence volumes; this chapter
reads the same model as an economist does, for its marginal effects, its
inference and the choice behaviour behind it.
\end{recall}

\begin{notation}[Symbols of this chapter]\label{not:b3:binary-outcomes:symbols}
$y_i \in \{0, 1\}$ is the outcome and $x_i$ the vector of regressors, with a
constant; $y_i^*$ is a \omterm{def:b3:binary-outcomes:logit-probit}{latent variable}. $\Lambda(z) = 1/(1 + \eu^{-z})$ is the
logistic and $\Phi$ the standard normal distribution function; $p_i$ is a
fitted probability. In the choice models, $U_{ij}$ is the utility of option
$j$ for person $i$ (a number, not a utility function of goods), $V_{ij}$ its
systematic part and $\eta_{ij}$ a taste shock. In demand, $\delta_j$ is a
product's mean utility (not depreciation) and $\alpha$ the price coefficient
(not a capital share).
\end{notation}

\section{The linear probability model}\label{sec:b3:binary-outcomes:lpm}

\begin{definition}[Linear probability model]\label{def:b3:binary-outcomes:lpm}
\index{linear probability model}
The \emph{linear probability model} is the regression of a binary outcome on
regressors, $y_i = x_i'\beta + u_i$, estimated by OLS. Since $\E[y_i \mid x_i] =
\Pr(y_i = 1 \mid x_i)$, its coefficients are changes in the probability.
\end{definition}

The model has real virtues. Its coefficients are read directly as effects on the
probability; it is the regression of the earlier chapters, so instruments, fixed
effects and differences in differences carry over unchanged; and as a linear
approximation to the conditional probability it is often good where most of the
data lie. Its errors are heteroskedastic by construction --- $\Var(u_i \mid
x_i) = p_i(1 - p_i)$ --- so robust standard errors are required. Its weakness is
at the edges: fitted values can fall below 0 or above 1, and a constant slope
is implausible for a probability near either.

\section{Logit and probit}\label{sec:b3:binary-outcomes:logit-probit}

\begin{definition}[Latent variable; probit model; logit model]\label{def:b3:binary-outcomes:logit-probit}
\index{latent variable}\index{probit model}\index{logit model}
Suppose $y_i = 1$ if a \emph{latent variable} $y_i^* = x_i'\beta + \eta_i$ is
positive and $y_i = 0$ otherwise, with $\eta_i$ independent of $x_i$ and
symmetric about zero with distribution function $F$. Then $\Pr(y_i = 1 \mid x_i)
= F(x_i'\beta)$. The \emph{probit model} takes $F = \Phi$, the standard normal;
the \emph{logit model} takes $F = \Lambda$, the logistic.
\end{definition}

The \omterm{def:b3:binary-outcomes:logit-probit}{latent variable} is a device, and sometimes an interpretation: the gain from
working, the net value of buying. Its scale is not identified --- multiplying
$\beta$ and the spread of $\eta$ by the same number changes nothing observed ---
so the variance of $\eta$ is fixed by convention: 1 for the probit, $\pi^2/3$ for
the logit (here $\pi$ is the circle's constant, not inflation). That is why logit
coefficients are larger than probit coefficients on the same data --- by a factor
between 1.80 and 1.95 in \cref{out:b3:binary-outcomes:three} below --- and why
neither set is comparable with the linear model's.

\begin{omfigure}
\begin{tikzpicture}[econ, xscale=1.1, yscale=2.2]
  \draw[econ axes] (-3.2,0) -- (3.4,0) node[right] {$\eta$};
  \draw[curve, domain=-3:3, samples=60, omDef] plot (\x, {exp(-\x*\x/2)/sqrt(2*pi)*2});
  \fill[omDef!20, domain=-0.8:3, samples=40] (-0.8,0) -- plot (\x, {exp(-\x*\x/2)/sqrt(2*pi)*2}) -- (3,0) -- cycle;
  \draw[guide] (-0.8,0) node[below] {$-x_i'\beta$} -- (-0.8,0.95);
  \node[font=\footnotesize, align=left, anchor=west] (lab) at (2.0,0.62) {$\Pr(y_i = 1)$\\$= \Pr(\eta_i > -x_i'\beta)$};
  \draw[->, thin] (lab.south west) ++(0.1,0) -- (1.3,0.18);
\end{tikzpicture}
\par\smallskip{\small The latent-variable picture. The outcome is 1 when the
shock $\eta_i$ exceeds $-x_i'\beta$; the shaded area is $F(x_i'\beta)$. A larger
$x_i'\beta$ moves the threshold left and raises the probability.}
\end{omfigure}

\section{The likelihood and its maximum}\label{sec:b3:binary-outcomes:likelihood}

\begin{definition}[Likelihood function; log-likelihood; maximum likelihood estimator]\label{def:b3:binary-outcomes:mle}
\index{likelihood function}\index{log-likelihood}\index{maximum likelihood estimator}
For independent observations with densities or probabilities $f(y_i \mid x_i;
\theta)$, the \emph{likelihood function} is $L(\theta) = \prod_i f(y_i \mid x_i;
\theta)$, the probability of the observed data as a function of the
parameters; the \emph{log-likelihood} is $\ell(\theta) = \ln L(\theta)$; and the
\emph{maximum likelihood estimator} $\hat\theta$ maximises it.
\end{definition}

For a binary model with $p_i = F(x_i'\beta)$, each observation contributes $p_i$
if $y_i = 1$ and $1 - p_i$ if not:
\begin{equation}\label{eq:b3:bo:loglik}
  \ell(\beta) = \sum_i \bigl[\, y_i \ln p_i + (1 - y_i) \ln(1 - p_i) \bigr].
\end{equation}

\begin{proposition}[The logit log-likelihood is concave; its score equations]\label{prop:b3:binary-outcomes:concave}
For the \omterm{def:b3:binary-outcomes:logit-probit}{logit model}, the gradient (the score) and the Hessian of $\ell$ are
\[
  \nabla \ell(\beta) = \sum_i (y_i - p_i)\, x_i, \qquad
  \nabla^2 \ell(\beta) = -\sum_i p_i (1 - p_i)\, x_i x_i'.
\]
The Hessian is negative semidefinite, so $\ell$ is concave, and any solution of
the score equations $\sum_i (y_i - p_i) x_i = 0$ is a global maximum. If $x_i$
contains a constant, the fitted probabilities average to the observed share of
ones; if it contains a dummy for a group, they do so within the group.
\end{proposition}

\begin{proof}
With $z_i = x_i'\beta$ and $p_i = \Lambda(z_i)$, $\Lambda'(z) = \Lambda(z)(1 -
\Lambda(z))$, so $\partial \ln p_i/\partial z_i = 1 - p_i$ and $\partial \ln(1 -
p_i)/\partial z_i = -p_i$. The derivative of observation $i$'s term with respect
to $z_i$ is $y_i(1 - p_i) - (1 - y_i)p_i = y_i - p_i$; the chain rule gives the
score. Differentiating again, $\partial (y_i - p_i)/\partial z_i = -p_i(1 - p_i)$,
which gives the Hessian. For any vector $v$, $v' \nabla^2\ell\, v = -\sum_i
p_i(1 - p_i)(x_i'v)^2 \le 0$. A concave function's stationary points are global
maxima. Finally, the score equation for the coefficient on a regressor that
equals 1 in a group and 0 elsewhere reads $\sum_{i \in \text{group}} (y_i - p_i) =
0$.
\end{proof}

\begin{proposition}[Properties of the maximum likelihood estimator]\label{prop:b3:binary-outcomes:mle}
If the model is correctly specified and regularity conditions hold, the \omterm{def:b3:binary-outcomes:mle}{maximum
likelihood estimator} is consistent and asymptotically normal, with variance
estimated by the inverse of minus the Hessian at $\hat\theta$ (the inverse of
the observed information).
\end{proposition}

\admitted
The proof expands the score around the true parameter and applies the law of
large numbers and the central limit theorem to its terms; it is in graduate
texts on econometric theory.

There is no formula for the maximum, so it is found by iteration. Newton's
method replaces the \omterm{def:b3:binary-outcomes:mle}{log-likelihood} near the current guess by its second-order
Taylor expansion and jumps to that quadratic's maximum: $\beta_{k+1} = \beta_k -
[\nabla^2\ell(\beta_k)]^{-1} \nabla\ell(\beta_k)$. The chart shows the simplest
case, a logit with only a constant $b$ on the opening chart's data, where the
maximum is the log of the odds of working, $\ln(13{,}820/2{,}172) = 1.8505$.
% figdata: newton

\begin{omchart}{The log-likelihood of a logit with only a constant $b$, on the
15{,}992 men of the opening chart, and Newton's steps from $b = 0$.}{Dehejia and
Wahba's NSW files (survey comparison sample); our computation}
\begin{tikzpicture}
\begin{axis}[omtime, width=0.7\linewidth, height=0.4\linewidth, xmin=-0.5, xmax=3.5,
  ymin=-16000, ymax=-5000, xtick={0,0.5,...,3}, xlabel={Constant $b$}, ylabel={Log-likelihood},
  x tick label style={/pgf/number format/.cd, fixed, precision=1},
  scaled y ticks=false, y tick label style={/pgf/number format/.cd, fixed, 1000 sep={\,}}]
  \addplot[omDef, thick, no marks, sharp plot] table[x=b, y=ll] {figdata/bachelor-3/05-binary-outcomes/out/loglik.csv};
  \addplot[omThm, only marks, mark=*, mark size=2pt,
    restrict expr to domain={\thisrow{step}}{0:1},
    nodes near coords={\pgfmathprintnumber[fixed, precision=0]{\pgfplotspointmeta}},
    point meta=\thisrow{step}, every node near coord/.append style={font=\scriptsize, anchor=south east}]
    table[x=b, y=ll] {figdata/bachelor-3/05-binary-outcomes/out/newton.csv};
  \addplot[omThm, only marks, mark=*, mark size=2pt, restrict expr to domain={\thisrow{step}}{2:9}]
    table[x=b, y=ll] {figdata/bachelor-3/05-binary-outcomes/out/newton.csv};
  \node[font=\scriptsize, anchor=south west, omThm] at (axis cs:1.9,-6400) {steps 2--4};
\end{axis}
\end{tikzpicture}
\end{omchart}

From $b = 0$ (every probability one half) the first step lands at 1.46, the
second at 1.80, the third at 1.850, and the fourth agrees with the maximum to
four decimals. % figdata: newton
\Cref{lst:b3:binary-outcomes:newton} runs the same method on the five-parameter
model of \cref{lst:b3:binary-outcomes:three}.

\omcode{code/bachelor-3/05-binary-outcomes/listings/lpm_logit_probit.py}{1}{25}{Three models of the same outcome, with \omterm{def:b3:binary-outcomes:ame}{average marginal effects}.}\label{lst:b3:binary-outcomes:three}

\omoutput{code/bachelor-3/05-binary-outcomes/outputs/lpm_logit_probit.txt}{Coefficients; then \omterm{def:b3:binary-outcomes:ame}{average marginal effects} of the logit and probit.}\label{out:b3:binary-outcomes:three}

\omcode{code/bachelor-3/05-binary-outcomes/listings/newton_logit.py}{1}{21}{Newton--Raphson by hand, with the library function \texttt{logit\_newton}.}\label{lst:b3:binary-outcomes:newton}

\omoutput{code/bachelor-3/05-binary-outcomes/outputs/newton_logit.txt}{The \omterm{def:b3:binary-outcomes:mle}{log-likelihood} climbs, then stops moving.}\label{out:b3:binary-outcomes:newton}

The \omterm{def:b3:binary-outcomes:mle}{log-likelihood} rises from $-11{,}084.8$ to $-5{,}147.84$ in five steps; the
last two change nothing that is printed, the quadratic convergence of
Newton's method. The hand-made estimates and standard errors agree with the
library's to every digit. % output: newton_logit

\section{Marginal effects and how to report them}\label{sec:b3:binary-outcomes:marginal}

A logit coefficient is the change in the log odds, $\ln[p/(1 - p)]$, for a unit
change in a regressor. Few readers think in log odds. The change in the
probability itself depends on where one starts: $\partial p_i/\partial x_{ik} =
F'(x_i'\beta)\,\beta_k$, which is largest when $p_i$ is one half and vanishes at
the edges.

\begin{definition}[Average marginal effect]\label{def:b3:binary-outcomes:ame}
\index{average marginal effect}
The \emph{average marginal effect} of a continuous regressor $x_k$ in a binary
model is the sample average of $F'(x_i'\hat\beta)\,\hat\beta_k$; for a dummy, the
average over the sample of the change in the fitted probability when the dummy
is switched from 0 to 1.
\end{definition}

In \cref{out:b3:binary-outcomes:three} the logit and probit coefficients differ
by a factor of about two, as their scales require, but their \omterm{def:b3:binary-outcomes:ame}{average marginal
effects} agree to within 0.002: a thousand dollars more of 1975 earnings goes with
a probability of working in 1978 higher by 1.35 points (logit) or 1.31 (probit).
The linear model's coefficient, 1.60 points, is larger: it is a single slope
for everyone, and fits the steep rise near zero earnings as well as the flat top.
And 15.7~\% of its fitted probabilities exceed one. % output: lpm_logit_probit

\begin{measure}[Reporting a logit: odds ratios and marginal effects]
Three numbers are reported for a logit coefficient, and they answer different
questions. The coefficient $\beta_k$ is the change in the log odds. Its
exponential $\eu^{\beta_k}$ is the odds ratio: the factor by which the odds
$p/(1 - p)$ are multiplied when $x_k$ rises by one, the same for everyone in the
model; it is the standard in medicine. The \omterm{def:b3:binary-outcomes:ame}{average marginal effect} is the
change in the probability, averaged over the people in the sample; it is the
standard in economics, because policy questions are about how many more
people will work, buy or move. The effect at the mean, computed at the average
values of the regressors, is a fourth number, often close to the average effect
and sometimes far from it, because the average person may resemble no one:
the average of a dummy for marriage is a person 0.7 married.
\end{measure}

\section{Discrete choice from random utility}\label{sec:b3:binary-outcomes:choice}

With more than two options --- car, bus, train, air; four brands of a good ---
the \omterm{def:b3:binary-outcomes:logit-probit}{latent variable} becomes a utility for each option.

\begin{definition}[Random utility model; multinomial logit]\label{def:b3:binary-outcomes:random-utility}
\index{random utility model}\index{multinomial logit}
In a \emph{random utility model}, person $i$ facing options $j = 1, \dots, J$
derives utility $U_{ij} = V_{ij} + \eta_{ij}$ from option $j$, where $V_{ij}$
depends on observed characteristics of the option and the person and
$\eta_{ij}$ is a taste shock, and chooses the option with the highest utility.
When the $\eta_{ij}$ are independent across options and people with the
extreme-value (Gumbel) distribution, $\Pr(\eta \le t) = \exp(-\eu^{-t})$, the
model is the \emph{multinomial logit} (or conditional logit when $V_{ij}$
depends on characteristics of the options).
\end{definition}

\begin{proposition}[Logit choice probabilities]\label{prop:b3:binary-outcomes:gumbel}
In the \omterm{def:b3:binary-outcomes:random-utility}{multinomial logit}, $\Pr(i \text{ chooses } j) = \eu^{V_{ij}} / \sum_k
\eu^{V_{ik}}$.
\end{proposition}

\begin{proof}
We prove it for two options; the general case follows by the same
integration over all options. Option 1 is chosen when $\eta_2 < V_1 - V_2 +
\eta_1$. Conditional on $\eta_1 = t$, this has probability $\exp(-\eu^{-(V_1 - V_2
+ t)})$. The density of $\eta_1$ is $\eu^{-t}\exp(-\eu^{-t})$. Writing $c = 1 +
\eu^{-(V_1 - V_2)}$, the probability is
\[
  \int_{-\infty}^{\infty} \eu^{-t} \exp(-c\,\eu^{-t})\, dt
  = \Bigl[\tfrac{1}{c}\exp(-c\,\eu^{-t})\Bigr]_{-\infty}^{\infty} = \frac1c
  = \frac{\eu^{V_1}}{\eu^{V_1} + \eu^{V_2}},
\]
substituting $s = \eu^{-t}$ in the integral if one prefers.
\end{proof}

\begin{definition}[Independence of irrelevant alternatives]\label{def:b3:binary-outcomes:iia}
\index{independence of irrelevant alternatives}
A choice model has the property of \emph{independence of irrelevant
alternatives} (IIA) if the ratio of the probabilities of any two options does not
depend on which other options are available or on their characteristics.
\end{definition}

\begin{proposition}[Logit has IIA]\label{prop:b3:binary-outcomes:iia}
In the \omterm{def:b3:binary-outcomes:random-utility}{multinomial logit}, $\Pr(j)/\Pr(k) = \eu^{V_{ij} - V_{ik}}$. Adding or
removing an option, or changing its characteristics, leaves the ratio unchanged;
a new option draws its users from the others in proportion to their
probabilities.
\end{proposition}

\begin{proof}
The denominators of $\Pr(j)$ and $\Pr(k)$ are the same sum and cancel. The
ratio involves only $V_{ij}$ and $V_{ik}$. If a new option takes probability
$q$, every other probability is multiplied by $1 - q$, since the ratios among
them are fixed and they sum to $1 - q$.
\end{proof}

The textbook warning is the red bus and the blue bus. If commuters split evenly
between a car and a red bus, and an identical blue bus is added, common sense
says the buses share the bus riders: car one half, each bus one quarter. Logit
says each option gets one third, because the new bus must take from the car in
proportion. Whenever options are closer substitutes for some than for others,
IIA misleads; nested logit, which groups similar options, and models whose
taste coefficients vary across people relax it.

\Cref{lst:b3:binary-outcomes:travel} estimates a conditional logit on real
data: 210 travellers, each choosing air, train, bus or car for one journey,
with each option's cost, time in the vehicle and time waiting at the terminal.

\omcode{code/bachelor-3/05-binary-outcomes/listings/travel_mode.py}{1}{22}{A conditional logit by Newton's method (\texttt{clogit}), and the forecast without air travel.}\label{lst:b3:binary-outcomes:travel}

\omoutput{code/bachelor-3/05-binary-outcomes/outputs/travel_mode.txt}{Coefficients with standard errors; car is the base mode.}\label{out:b3:binary-outcomes:travel}

The ratio of the time and cost coefficients is a price of time: travellers
trade an hour in the vehicle for 17.2 dollars. % output: travel_mode
The model fits the four shares exactly, a consequence of the mode constants:
the multinomial analogue of the score equations in
\cref{prop:b3:binary-outcomes:concave}. Without air travel, it sends the air
travellers to the other modes in proportion, for each traveller, to their
probabilities. In these data the households of air travellers have incomes
like those of car travellers and well above those of bus travellers; common
sense would send many of them to the car and few to the bus, and IIA cannot.

\begin{record}[Forecasting a new rail line from a logit model]
Before the Bay Area Rapid Transit system opened in San Francisco, McFadden's
group surveyed commuters in 1972, estimated a \omterm{def:b3:binary-outcomes:random-utility}{multinomial logit} of their
journey to work, and used it to forecast the choices of 631 commuters surveyed
again in 1975, after the new line opened. The forecast share of work trips by
the new line was 6.3~\% (standard error 2.5~\%); the realised share was
6.2~\%. % ledger: b3bo:mcf.bart
The forecasts for driving alone (55.8~\% against 59.9~\%) and for the bus
(14.9~\% against 12.2~\%) were less exact, and McFadden notes that the model
overestimated people's willingness to walk to the new stations.
\end{record}

\subsection*{Logit demand for products}

The same model, applied to the market shares of differentiated products, is the
workhorse of empirical industrial organisation. Each consumer buys one of $J$
products or nothing (option 0, with $V_0 = 0$), and $V_j = \delta_j = x_j'\gamma
+ \alpha p_j + \xi_j$, where $\xi_j$ is quality seen by consumers but not by the
economist.

\begin{definition}[Log-share inversion]\label{def:b3:binary-outcomes:inversion}
\index{log-share inversion}
In the logit demand model, the \emph{log-share inversion} recovers each
product's mean utility from its market share and the outside share: $\ln s_j -
\ln s_0 = \delta_j = x_j'\gamma + \alpha p_j + \xi_j$, a linear equation that can be
estimated by regression across products and markets (Berry 1994).
\end{definition}

The inversion follows from $s_j/s_0 = \eu^{\delta_j}/\eu^{0}$. Its value is
that it turns a nonlinear choice model into a regression, so that the tools of
\cref{ch:b3:instrumental-variables} apply: prices are set by firms that know
$\xi_j$, so $p_j$ is correlated with the error, and OLS is biased toward zero
(\cref{exo:b3:binary-outcomes:11}). Instruments are costs, or characteristics
of rival products.

\omcode{code/bachelor-3/05-binary-outcomes/listings/logit_demand.py}{1}{27}{Logit demand on 300 simulated markets: the inversion and the elasticities.}\label{lst:b3:binary-outcomes:demand}

\omoutput{code/bachelor-3/05-binary-outcomes/outputs/logit_demand.txt}{Own and cross price elasticities in one town.}\label{out:b3:binary-outcomes:demand}

The own elasticity is $\alpha p_j(1 - s_j)$ and the cross elasticity of $j$'s
share with respect to $k$'s price is $-\alpha p_k s_k$, the same for every $j$:
in each column of \cref{out:b3:binary-outcomes:demand}, a rise in product
$k$'s price raises every rival's share by the same percentage. That is IIA again,
and it is why logit demand, used to judge a merger, says that two products
that are close substitutes compete no harder than two that are not.

\begin{model}{Random utility and logit demand}
\assumes Each person chooses the option with the highest utility, a mean
utility depending on observed characteristics plus a taste shock independent
across options and people with the extreme-value distribution.
\predicts Choice probabilities $\eu^{V_j}/\sum_k \eu^{V_k}$; fitted shares equal to
observed shares when each option has a constant; substitution in proportion to
shares (IIA); a price of time equal to the ratio of time and cost coefficients.
\tested In forecasting the share of a new rail line in 1975 from 1972 data, the
logit came within 0.1 percentage point; its forecasts for other modes were
off by up to four points.
\breaks When options are close substitutes for each other (two buses, two
similar cars), IIA sends too many customers to unlike options; nested logit and
random-coefficient models relax it. When prices respond to unobserved quality,
estimates of the price coefficient need instruments.
\end{model}

\section{Exercises}\label{sec:b3:binary-outcomes:exercises}

\begin{exercise}[$\star$]\label{exo:b3:binary-outcomes:1}
In the opening chart, where do the three fitted lines differ most, and where do
all three fail? Which feature of the data explains the failure?
\end{exercise}

\begin{exercise}[$\star$]\label{exo:b3:binary-outcomes:2}
A logit has $x'\hat\beta = 0.8$ for a person, and the coefficient on a regressor
is 0.5. Compute the person's probability and the marginal effect of the
regressor on it. Repeat for $x'\hat\beta = 3$.
\end{exercise}

\begin{exercise}[$\star$]\label{exo:b3:binary-outcomes:3}
The logit coefficient on being married in \cref{out:b3:binary-outcomes:three}
is $-0.2981$. Compute the odds ratio and say what it means.
\end{exercise}

\begin{exercise}[$\star$]\label{exo:b3:binary-outcomes:4}
Describe data on commuters' choices before and after a new tram line opened
that would contradict IIA.
\end{exercise}

\begin{exercise}[$\star\star$]\label{exo:b3:binary-outcomes:5}
Add a dummy for no earnings in 1975 to the logit of the opening chart.
\begin{enumerate}
\item Compare the \omterm{def:b3:binary-outcomes:mle}{log-likelihoods} and the mean fitted probability for men with no
earnings in 1975 with the observed share.
\item With controls for age, education and marriage, report the \omterm{def:b3:binary-outcomes:ame}{average marginal
effect} of the dummy and of 1975 earnings.
\end{enumerate}
\end{exercise}

\begin{exercise}[$\star\star$]\label{exo:b3:binary-outcomes:6}
Five observations have $y = (1, 0, 1, 1, 0)$ and fitted logit probabilities
$p = (0.8, 0.3, 0.6, 0.9, 0.5)$. Compute the \omterm{def:b3:binary-outcomes:mle}{log-likelihood}. Which observation
contributes the most to it in absolute value, and why?
\end{exercise}

\begin{exercise}[$\star\star$]\label{exo:b3:binary-outcomes:7}
Using the elasticities printed in \cref{out:b3:binary-outcomes:demand}, by how
many percent do the shares of products 2 and 3 change if product 1's price rises
by 1~\%? And if product 4's price rises by 1~\%? Why are the answers the same for
products 2 and 3?
\end{exercise}

\begin{exercise}[$\star\star$]\label{exo:b3:binary-outcomes:8}
A logit with only a constant is fitted to $n$ observations with $n_1$ ones.
\begin{enumerate}
\item Write the \omterm{def:b3:binary-outcomes:mle}{log-likelihood} and the Newton step from $b$.
\item With $n = 100$ and $n_1 = 80$, compute two Newton steps from $b = 0$ and
compare with the maximum, $\ln 4$.
\end{enumerate}
\end{exercise}

\begin{exercise}[$\star\star$]\label{exo:b3:binary-outcomes:9}
How would you check, in the \omterm{def:b3:binary-outcomes:lpm}{linear probability model} of
\cref{lst:b3:binary-outcomes:three}, that the linear form is wrong? Give two
checks that use only its fitted values and the data.
\end{exercise}

\begin{exercise}[$\star\star\star$]\label{exo:b3:binary-outcomes:10}
Show that in a logit with a constant, the fitted probabilities average to the
observed share, and that in a logit whose only regressors are dummies for
mutually exclusive groups (and no constant), the maximum likelihood fitted
probability in each group is the group's observed share. What does that say
about the opening chart's first dot in \cref{exo:b3:binary-outcomes:5}?
\end{exercise}

\begin{exercise}[$\star\star\star$]\label{exo:b3:binary-outcomes:11}
Repeat the simulation of \cref{lst:b3:binary-outcomes:demand}, but let firms
raise the price of products with high unobserved quality: $p_j = 1 + 0.3 x_j +
c_j + 0.8\,\xi_j$, where $c_j$ is a cost that does not enter demand.
\begin{enumerate}
\item Estimate the price coefficient by OLS on the \omterm{def:b3:binary-outcomes:inversion}{log-share inversion}. In which
direction is it biased, and why?
\item Estimate it by \omterm{def:b3:instrumental-variables:tsls}{two-stage least squares} with $c_j$ as the instrument.
\end{enumerate}
\end{exercise}

\begin{exercise}[$\star\star\star$]\label{exo:b3:binary-outcomes:12}
In a town, 60~\% of commuters drive, 30~\% take the bus and 10~\% cycle. A new
cycle lane will make cycling attractive enough to take 10~\% of all commuters
from the other modes (a new option, ``cycle lane'').
\begin{enumerate}
\item Under IIA, what are the new shares?
\item Why is that forecast doubtful, and in which direction is it likely to err
for each mode?
\end{enumerate}
\end{exercise}

\section{Problem: A New Bus Line in Brenmoor}\label{sec:b3:binary-outcomes:problem}

\begin{problem}[{Weekend problem --- a new bus line in Brenmoor}]\label{pb:b3:binary-outcomes:bus}
Brenmoor (invented) surveys 1{,}500 commuters. Each commutes by car, bus or
bicycle; the survey records, for each commuter and each mode, the cost of the
trip in dalens (the bus fare is \dal{2.5}, cycling costs nothing) and the
door-to-door time in minutes. Of the 1{,}500, 992 chose the car or the bus.
% figdata: commuters_meta
A statistician estimates two models by maximum likelihood. The first is a binary
logit of choosing the car rather than the bus, among those 992, on the
differences car minus bus in cost and in time. The second is a three-way
conditional logit with a constant for the car and for the bicycle (the bus is
the base), a cost and a time coefficient. % figdata: commuters_binary, commuters_mnl

\begin{center}\small
\begin{minipage}[t]{0.47\linewidth}\centering
Binary logit, car against bus\par\smallskip
\pgfplotstabletypeset[col sep=comma,
  columns/term/.style={column name={Term}, string type, column type=l},
  columns/est/.style={column name={Estimate}, fixed, precision=4, zerofill},
  columns/se/.style={column name={s.e.}, fixed, precision=4, zerofill},
  every head row/.style={after row=\hline}, every last row/.style={after row=\hline}]
  {figdata/bachelor-3/05-binary-outcomes/out/commuters_binary.csv}
\end{minipage}\hfill
\begin{minipage}[t]{0.47\linewidth}\centering
Three-way logit\par\smallskip
\pgfplotstabletypeset[col sep=comma,
  columns/term/.style={column name={Term}, string type, column type=l},
  columns/est/.style={column name={Estimate}, fixed, precision=4, zerofill},
  columns/se/.style={column name={s.e.}, fixed, precision=4, zerofill},
  every head row/.style={after row=\hline}, every last row/.style={after row=\hline}]
  {figdata/bachelor-3/05-binary-outcomes/out/commuters_mnl.csv}
\end{minipage}

\medskip
\pgfplotstabletypeset[col sep=comma,
  columns/mode/.style={column name={Mode}, string type, column type=l},
  columns/actual/.style={column name={Observed share}, fixed, precision=3, zerofill, empty cells with={---}},
  columns/fitted/.style={column name={Fitted, three-way logit}, fixed, precision=3, zerofill, empty cells with={---}},
  columns/forecast/.style={column name={Forecast with the express}, fixed, precision=3, zerofill},
  every head row/.style={after row=\hline}, every last row/.style={after row=\hline}]
  {figdata/bachelor-3/05-binary-outcomes/out/commuters_shares.csv}
\end{center}
The forecast adds to the three-way model an express bus with the ordinary bus
fare, the bus constant and a door-to-door time 10 minutes shorter than the bus
for every commuter.

\medskip\noindent\textbf{Part I --- car or bus.}
\begin{enumerate}
\item Write the probability of choosing the car in the binary logit, in terms of
the cost and time differences.
\item A commuter's car trip costs \dal{4.5} and takes 30 minutes; the bus costs
\dal{2.5} and takes 45 minutes. Compute $x'\hat\beta$ and the probability that she
drives.
\item Interpret the constant. Why does it matter that the bus is the base?
\item Is the cost coefficient significantly different from zero? Give its
$t$-statistic.
\item Why does the logit need no assumption about the variance of the taste
shocks to be estimated, but fix it anyway?
\item Why might estimating on the 992 car and bus users only, leaving out the
cyclists, still give consistent estimates in this model?
\end{enumerate}

\medskip\noindent\textbf{Part II --- marginal effects and the value of time.}
\begin{enumerate}[resume]
\item For the commuter of Question 2, compute the marginal effect of the bus
fare on her probability of driving.
\item Compute the effect of cutting the bus time by 10 minutes on the same
probability, using the marginal effect, and then exactly. Why do they differ?
\item Compute the value of travel time from the binary logit, in dalens per hour.
\item Compute it from the three-way logit.
\item Give one reason why the two estimates differ.
\item A road scheme would save each of 20{,}000 car commuters 5 minutes a day
for 220 days a year. Value the saving with the three-way estimate.
\end{enumerate}

\medskip\noindent\textbf{Part III --- three modes.}
\begin{enumerate}[resume]
\item Write the probability of cycling in the three-way logit.
\item Check that the fitted shares equal the observed shares, and explain why.
\item Interpret the bicycle constant.
\item For a commuter with $V_{\text{car}} = -2.0$, $V_{\text{bus}} = -3.0$ and
$V_{\text{bicycle}} = -2.5$, compute the three probabilities.
\item For the same commuter, compute the ratio of the car and bus probabilities.
What happens to it if the bicycle becomes unavailable?
\item Is the time coefficient the same for the three modes? Is that a
reasonable assumption for cycling?
\end{enumerate}

\medskip\noindent\textbf{Part IV --- the express line.}
\begin{enumerate}[resume]
\item What share does the model forecast for the express?
\item From which mode does it take the most riders, in percentage points, and in
percent of each mode's share?
\item Why does the model take riders from each mode in nearly equal proportions?
\item Which mode would you expect the express to take most of its riders from in
fact, and why?
\item How would a nested logit, with the two buses in one nest, change the
forecast?
\item What survey, run after the express opens, would test the forecast, as
McFadden's group did for a new rail line?
\item State the results: the commuters' value of travel time, in dalens per hour,
from the three-way model, and the express line's forecast share, with the
caveat that applies to it.
\end{enumerate}
\end{problem}
